\documentclass[11pt,a4paper]{article}

% ------------------------------------------------------------
% Packages
% ------------------------------------------------------------
\usepackage[utf8]{inputenc}
\usepackage[T1]{fontenc}
\usepackage{lmodern}
\usepackage[english]{babel}
\usepackage[a4paper,margin=2.5cm]{geometry}
\usepackage{amsmath,amssymb}
\usepackage{siunitx}
\usepackage{booktabs}
\usepackage{array}
\usepackage{graphicx}
\usepackage{float}
\usepackage{caption}
\usepackage{subcaption}
\usepackage{pgfplots}
\usepackage{pgfplotstable}
\usepackage{xcolor}
\usepackage{enumitem}
\usepackage{hyperref}
\usepackage{microtype}
% \usepackage{siunitx}

% ------------------------------------------------------------
% Custom commands
% ------------------------------------------------------------
\renewcommand{\thesubsection}{\alph{subsection})}
\renewcommand{\thesection}{}
\newcommand{\degCel}{\,^\circ\mathrm{C}}
\newcommand{\Frac}[2]{\left(\frac{#1}{#2}\right)}




\pgfplotsset{compat=1.18}

\sisetup{
    per-mode=symbol,
    exponent-product=\cdot,
    output-decimal-marker={.},
    group-separator={,},
    group-minimum-digits=4
}

\hypersetup{
    colorlinks=true,
    linkcolor=black,
    citecolor=black,
    urlcolor=blue,
    pdftitle={Assignment 5 - Generator Design Tool}
}

% ------------------------------------------------------------
% Document information
% ------------------------------------------------------------
\title{
    \textbf{FENT2321}\\[0.3cm]
    Wind Energy and Design of a Wind Turbine\\[0.5cm]
    \Large Assignment 5: Generator Design Tool
}

\author{Markus Arntzen}
\date{}

\begin{document}

\maketitle
\thispagestyle{empty}

\clearpage
\tableofcontents
\clearpage

% ============================================================
\section{Problem 1: Induced Voltage}
% ============================================================

% ------------------------------------------------------------
% 1.1 Blade geometry
% ------------------------------------------------------------
\subsection{}
\begin{equation}
    \omega_m = 2\pi f = \frac{\pi \cdot RPM}{30} = \frac{\pi \cdot 1500}{30} = 157.1\frac{rad}{s}
\end{equation}

\subsection{}
\begin{equation}
    v_{lin} = \omega_m R = \omega_m \cdot (r + \frac{t_{coil}}{2}) =
    157.1 \times \left(0.045 + \frac{0.010}{2}\right) = 7.85 \frac{m}{s}
\end{equation}

\subsection{}
The loop length (longest part of the coil) laying perpendicular to the magnetic field lines
will induce a voltage. The total length will be $2\times$ the length of one side, giving;
$2\times l = 2\times 0.1m = 0.2m$.

\subsection{}
\begin{equation}
    V = NBv(2l) = 50\cdot0.2\cdot7.85\cdot(2\cdot0.1) = 15.7V
\end{equation}

\subsection{}
\begin{equation}
    V = \frac{NBv(2l)N_sf_w}{\sqrt{2}} = \frac{50\cdot0.2\cdot7.85\cdot(2\cdot0.1)\cdot2\cdot0.9}{\sqrt{2}} = 20.0V
\end{equation}

\vspace{0.5cm}

\section{Problem 2: Resistance}
\subsection{}
\begin{equation}
    L_{wind} = \left(2\cdot l + 2\cdot \frac{2\pi \left(r + \frac{t_{coil}}{2}\right)}{N_s}\right)\cdot L_{corr}
    = \left(2\cdot0.1 + 2\cdot \frac{2\pi \left(0.045 + \frac{0.01}{2}\right)}{2}\right)\cdot 1.10
    = 0.566m
\end{equation}

\subsection{}
\begin{equation}
    L_{Cu} = L_{turn}N_{turn}N_sN_{phase} = 0.566\cdot50\cdot2\cdot1 = 56.6m
\end{equation}

\subsection{}
\begin{equation}
    R_{Cu} = \frac{\rho_{Cu}(1 + \alpha_{Cu}(T-T_0))}{A_{Cu}}
    = \frac{1.68\cdot10^{-8}(1 + 0.0039(30\degCel - 20\degCel))}{7.85\cdot10^{-7}}
    = 0.022\frac{\Omega}{m}
\end{equation}

\subsection{}
\begin{equation}
     R_{phase} = R_{Cu}L_{stator}c_c = 0.022\cdot56.6\cdot1.10 = 1.38\Omega
\end{equation}

\vspace{0.5cm}

\section{Problem 3: Reactance and current}
\subsection{}
\begin{equation}
    L_s = L_{s0}\left(\frac{N_{turn}}{N_0}\right)^2N_s
    = 0.0008\cdot\left(\frac{50}{28}\right)^2\cdot2 = 5.10\cdot10^{-3}\unit{H}
\end{equation}

\subsection{}
\begin{equation}
    \omega_e = \omega_m\frac{N_{poles}}{2} = 157.1\frac{2}{2} = 157.1\frac{rad}{s}
\end{equation}

\subsection{}
\begin{equation}
    X_s = \omega_e L_s = 157.1\cdot5.10\cdot10^{-3} = 0.801\Omega
\end{equation}

\subsection{}
\begin{equation}
    |I| = \frac{V_{ind}}{\sqrt{(R_s + R_L)^2 + X_s^2}}
    = \frac{20}{\sqrt{(1.38+4.00)^2 + 0.801^2}}
    = 3.67\unit{A}
\end{equation}

\subsection{}
\begin{equation}
    V_{term} = IR_L = 3.67\cdot4.00 = 14.7\unit{V}
\end{equation}

\vspace{0.5cm}

\section{Problem 4: Power}
\subsection{}
\begin{equation}
    P_{load} = N_{phase}I^2R_L = 1\cdot3.67^2\cdot4.00 = 54\unit{W}
\end{equation}

\subsection{}
\begin{equation}
    P_{loss_{Cu}} = N_{phase}I^2R_s = 1\cdot3.67^2\cdot1.38 = 18.7\unit{W}
\end{equation}

\subsection{}
\begin{equation}
    P_{loss_m} = \Frac{RPM}{RPM_0}P_{loss_0} + F_{loss\,\,offset}
    = \Frac{1500}{1000}\cdot75 + 5 = 117.5\unit{W}
\end{equation}

\subsection{}
\begin{equation}
    P_{gen} = P_{load} + P_{loss_{Cu}} + P_{loss_m}
    = 54 + 18.7 + 117.5 = 190.1\unit{W}
\end{equation}

\subsection{}
\begin{equation}
    \eta = \Frac{P_{out}}{P_{in}}100 = \Frac{54}{190.1}\cdot100 = 28.4\%
\end{equation}

\subsection{}
\begin{equation}
    T_{gen} = \Frac{P_{in}}{\omega_m} = \Frac{190.1}{157.1} = 1.211\unit{Nm}
\end{equation}

\subsection{}
from assignment 4, the wind turbine torque came out as 1.2\,Nm which gives
$T_{diff} = 0.011$ with $T_{gen} > T_{aero}$. This means that at this rotational speed
$\omega_m$ will the generator brake the rotors speed until the torques become balanced.


\end{document}